%++++++++++++++++++++++++++++++++++++++++
\documentclass[letterpaper,12pt]{article}

\usepackage[margin=1in,letterpaper]{geometry}
\usepackage{amsmath}
\usepackage{graphicx}
\usepackage{siunitx}
\usepackage{cite}
\usepackage[final]{hyperref}

\hypersetup{
    colorlinks=true,
    linkcolor=blue,
    citecolor=blue,
    urlcolor=blue
}

\sisetup{separate-uncertainty=true}
%++++++++++++++++++++++++++++++++++++++++

\begin{document}

\title{Early Giants in a Young Universe: Mapping Black-Hole Growth Scenarios for JWST Accretion Candidates}
\author{Nicole Jiang}
\date{\today}
\maketitle

\begin{abstract}
\textit{Aims.} The James Webb Space Telescope (JWST) pushes sensitive spectroscopy and imaging into the first few billion years of cosmic history, where there is not much time for black holes and galaxies to assemble. That makes many high-redshift discoveries natural stress tests for models of early black-hole seeding, accretion, and host-galaxy assembly. However, key quantities like $M_{\rm BH}$, $M_{*}$, or $L_{\rm bol}$ are often inferred with different methods and assumptions across papers, even though small shifts in these quantities can imply very different growth histories.

\textit{Methods.} This project builds a standardized catalogue of cross-paper JWST-identified accretion candidates at $z \geq 4$ that tracks how quantities were inferred and under which assumptions. It then tests these objects against different seed and growth scenarios using the analytic black-hole growth model from Dayal (2024). We visualize whether an object can be produced under certain physical conditions using parameter-space maps in seed mass, Eddington fraction, seed redshift, and radiative efficiency.

\textbf{Keywords:} high-redshift galaxies; black holes; active galactic nuclei; JWST; accretion; black-hole seeds
\end{abstract}

\section{Introduction and Background}

The James Webb Space Telescope has pushed sensitive spectroscopy and imaging into the first few billion years of cosmic history. This is a regime where there is not much time for galaxies and black holes to assemble, so every massive or actively accreting black-hole candidate becomes a useful stress test. If a black hole is already large at high redshift, then either it started from a sufficiently massive seed, accreted efficiently for a long enough period of time, formed very early, or some combination of these conditions occurred.

This makes high-redshift accreting systems especially useful for studying early structure formation. They connect the observational side of JWST discoveries to theoretical questions about how the first black-hole seeds formed and how quickly they could grow. Light seeds, such as remnants of Population III stars, are physically natural but may require very strong or sustained accretion to reach high masses by early cosmic times. Heavy seeds reduce the required amount of later growth, but usually require more special formation environments. More speculative possibilities, such as primordial black-hole seeding, give the black hole an even earlier head start.

The difficulty is that the objects being compared across the literature are not always inferred in the same way. Quantities such as $M_{\rm BH}$, $M_{*}$, $L_{\rm bol}$, and $f_{\rm Edd}$ can depend strongly on the adopted method. Black-hole masses may come from virial estimators, line widths, luminosity conversions, or model-dependent assumptions. Stellar masses can shift depending on stellar-population modelling, dust assumptions, and how the AGN contribution is separated from the host galaxy. Because black-hole growth is exponential, even modest changes in these inferred quantities can move an object into or out of the ``peculiar'' regime.

This motivates the main purpose of the project: to build a standardized, source-tracked catalogue of JWST-identified accretion candidates at $z \geq 4$. The goal is not just to list objects, but to preserve where each measurement came from and what assumptions were used to infer it. This makes it easier to ask which objects are genuinely difficult to explain, rather than which objects only appear extreme because of inconsistent conventions across papers.

The project then uses the catalogue to test different seed and growth scenarios. For each observed black-hole mass at redshift $z$, the model asks what combinations of seed mass, average Eddington fraction, seed redshift, and radiative efficiency can reproduce the observed object. This produces parameter-space maps that show where each object lies relative to physically interpretable formation scenarios. In this way, the project acts as a high-redshift accretion atlas: a reproducible way to compare early black-hole candidates under common assumptions.

\section{Background Theory}

The available growth time is set by the age of the Universe at the observed redshift and at the assumed seed redshift. For a flat $\Lambda$CDM cosmology, the cosmic time in Gyr can be written as

\begin{equation}
t(z) =
\frac{2}{3H_0\sqrt{\Omega_{\Lambda}}}
\operatorname{asinh}
\left[
\frac{\sqrt{\Omega_{\Lambda}/\Omega_m}}{(1+z)^{3/2}}
\right],
\label{eq:cosmic_time}
\end{equation}

where $H_0$ is the Hubble constant, $\Omega_m$ is the matter density parameter, and $\Omega_{\Lambda}$ is the dark-energy density parameter. For each object, the time available for growth is the difference between the cosmic time at the observed redshift and the cosmic time at the seed redshift.

Following Dayal (2024), the black-hole growth model is

\begin{equation}
M_{\rm BH}(t) =
M_{\rm seed}(t_{\rm seed})
\exp
\left[
\frac{4\pi G m_p}{c\sigma_T}
\frac{1-\epsilon}{\epsilon}
f_{\rm Edd}
(t-t_{\rm seed})
\right],
\label{eq:bh_mass_growth}
\end{equation}

where $M_{\rm BH}(t)$ is the black-hole mass at time $t$, $M_{\rm seed}(t_{\rm seed})$ is the seed mass at the time the seed formed, $f_{\rm Edd}$ is the average Eddington fraction, and $\epsilon$ is the radiative efficiency. The constants $G$, $m_p$, $c$, and $\sigma_T$ are the gravitational constant, proton mass, speed of light, and Thomson scattering cross-section.

This equation shows why high-redshift black holes are so sensitive to assumptions. Growth depends exponentially on the available time, the average Eddington fraction, and the radiative efficiency term $(1-\epsilon)/\epsilon$. A higher seed mass gives the black hole a larger starting point. A higher $f_{\rm Edd}$ increases the growth rate. A lower radiative efficiency allows more of the accreted mass to remain in the black hole. A higher seed redshift gives the system more time to grow before it is observed.

The parameter maps in this project are built around these dependencies. The x-axis will represent $\log_{10}(M_{\rm seed}/M_{\odot})$, the y-axis will represent $f_{\rm Edd}$, and the colour scale will show the predicted $\log_{10}(M_{\rm BH}/M_{\odot})$ for a given choice of $z_{\rm seed}$ and $\epsilon$. Observed catalogue objects can then be placed on these maps to show which seed and growth assumptions are sufficient to reproduce them.

This approach helps separate different levels of tension. Some objects may be explained by ordinary seed masses and sub-Eddington or near-Eddington accretion. Others may require heavier seeds, earlier formation, lower radiative efficiencies, or sustained high accretion fractions. The most interesting candidates are the ones that remain difficult to reproduce even when reasonable alternative assumptions are allowed.

% Future sections can be added after the catalogue and plots are complete:
% \section{Methods}
% \section{Results}
% \section{Discussion}
% \section{Conclusion}

\begin{thebibliography}{99}

\bibitem{Dayal2024}
Dayal, P. 2024,
\textit{Astronomy \& Astrophysics}, 690, A182.
\href{https://www.aanda.org/articles/aa/full_html/2024/10/aa51481-24/aa51481-24.html}{A\&A article}

\end{thebibliography}

\appendix

\section{Derivation of the Cosmic Time Relation}

Start from flat $\Lambda$CDM:
\[
t(z)=\frac{1}{H_0}\int_z^\infty
\frac{dz'}{(1+z')\sqrt{\Omega_m(1+z')^3+\Omega_\Lambda}}.
\]

Let
\[
a\equiv \frac{1}{1+z},\qquad a'\equiv \frac{1}{1+z'},\qquad
dz'=-\frac{da'}{a'^2}.
\]

Then
\[
t(a)=\frac{1}{H_0}\int_0^a
\frac{da'}{a'\sqrt{\Omega_m/a'^3+\Omega_\Lambda}}
=
\frac{1}{H_0}\int_0^a
\frac{\sqrt{a'}\,da'}{\sqrt{\Omega_m+\Omega_\Lambda a'^3}}.
\]

Let
\[
u\equiv \sqrt{\frac{\Omega_\Lambda}{\Omega_m}}\,a'^{3/2},
\]
so that
\[
du=\sqrt{\frac{\Omega_\Lambda}{\Omega_m}}\frac{3}{2}a'^{1/2}da',
\]
and therefore
\[
a'^{1/2}da'=\frac{2}{3}\sqrt{\frac{\Omega_m}{\Omega_\Lambda}}\,du.
\]

Since
\[
\sqrt{\Omega_m+\Omega_\Lambda a'^3}
=
\sqrt{\Omega_m}\sqrt{1+u^2},
\]
we obtain
\[
t(a)=
\frac{1}{H_0}
\frac{2}{3\sqrt{\Omega_\Lambda}}
\int_0^{u(a)}
\frac{du}{\sqrt{1+u^2}}.
\]

Thus
\[
t(a)=
\frac{2}{3H_0\sqrt{\Omega_\Lambda}}
\operatorname{asinh}\!\big(u(a)\big),
\]
where
\[
u(a)=
\sqrt{\frac{\Omega_\Lambda}{\Omega_m}}\,a^{3/2}.
\]

Using $a=(1+z)^{-1}$ gives
\[
\boxed{
t(z)=
\frac{2}{3H_0\sqrt{\Omega_\Lambda}}
\operatorname{asinh}
\left[
\frac{\sqrt{\Omega_\Lambda/\Omega_m}}{(1+z)^{3/2}}
\right]
}.
\]

\end{document}